<?php
// save.php - Speichert die Umfragedaten in der Datenbank

// Datenbankverbindung einbinden
require __DIR__ . '/inc/db.php';

// ============================================================
// 1. POST-Daten auslesen
// ============================================================
$userId = isset($_POST['userId']) ? (int)$_POST['userId'] : 0;
$q1     = isset($_POST['Q1']) ? trim($_POST['Q1']) : '';
$q2     = isset($_POST['Q2']) ? trim($_POST['Q2']) : '';
$q3     = isset($_POST['Q3']) ? trim($_POST['Q3']) : '';
$q4     = isset($_POST['Q4']) ? trim($_POST['Q4']) : '';
$q5     = isset($_POST['Q5']) ? trim($_POST['Q5']) : '';

// ============================================================
// 2. Validierung: UserId muss vorhanden sein
// ============================================================
if ($userId === 0) {
    die("Fehler: Kein Nutzer gefunden. Bitte die Umfrage erneut starten.");
}

// ============================================================
// 3. Aktuelles Datum mit date() erzeugen
// ============================================================
$datum = date('Y-m-d');  // Format: Jahr-Monat-Tag (z.B. 2026-07-25)

// ============================================================
// 4. Daten in survey_table speichern
// ============================================================
try {
    $stmt = $pdo->prepare("
        INSERT INTO survey_table (UserId, Datum, Q1, Q2, Q3, Q4, Q5)
        VALUES (?, ?, ?, ?, ?, ?, ?)
    ");
    
    $stmt->execute([$userId, $datum, $q1, $q2, $q3, $q4, $q5]);
    
    $success = true;
} catch (PDOException $e) {
    $error = "Speicherung fehlgeschlagen: " . $e->getMessage();
    $success = false;
}

// ============================================================
// 5. Ergebnis anzeigen
// ============================================================
?>
<!DOCTYPE html>
<html lang="de">
<head>
    <meta charset="UTF-8">
    <title>Umfrage gespeichert</title>
</head>
<body>
    <h1>Umfrage zum Studium</h1>

    <?php if ($success): ?>
        <hr>
        <p style="font-size: 1.2em; font-weight: bold; text-align: center;">
            Die Umfrage wurde erfolgreich gespeichert!
        </p>
        <hr>
        <p style="text-align: center;">
            <a href="report.php">Report</a>
        </p>
    <?php else: ?>
        <p style="color: red; font-weight: bold;">
            <?php echo htmlspecialchars($error, ENT_QUOTES, 'UTF-8'); ?>
        </p>
        <p><a href="index.php">Zurück zur Startseite</a></p>
    <?php endif; ?>

</body>
</html>